\chapter{الشغل والطاقة الميكانيكية}\label{ch:g12:work-and-energy}

أسقِط كرة زجاجية في سلطانية: فتغوص، وتتجاوز القاع، و
تصعد إلى ارتفاع انطلاقها، وتتردد، ثم ترتدّ --- ويمكن التنبؤ بكل
ارتداد من منحنٍ واحد (هو مقطع السلطانية) وخط أفقي
واحد (هو \omterm{def:g10:energy-conservation:energy}{طاقة} الكرة)، دون حلّ أي معادلة حركة.
وقد سعّرت السنة الماضية القوى على المستقيمات (\cref{ch:g11:work-of-force})
ووازنت \omterm{def:g11:mechanical-energy:kinetic}{الطاقة الحركية} بالكامنة
(\cref{ch:g11:mechanical-energy})؛ وهذه السنة تذهب المحاسبة
إلى كل مكان: المسارات المنحنية، والزنبركات، والمناظر بنظرة واحدة.

\section{الشغل على أي مسار}

\begin{definition}[الشغل على مسار]\label{def:g12:work-and-energy:path}
لتؤثر \omterm{def:g10:inertia:force}{قوة} $\vect F$ في جسم يتحرك على أي \omterm{def:g10:relative-motion:trajectory}{مسار} من $A$ إلى
$B$: قطّع \omterm{def:g10:relative-motion:trajectory}{المسار} إلى انتقالات $\vect{\dd\ell}$ قصيرة إلى حدّ يجعل
كلًّا منها مستقيمًا ولا تكاد $\vect F$ تتغير عليه. فتكسب كل خطوة
\omterm{def:g11:work-of-force:work}{الشغل} الصغير $\scal{F}{\dd\ell} = F\,\dd\ell\cos\theta$
(\cref{ch:g11:work-of-force})، ويكون
\emph{الشغل على المسار}\index{شغل (قوة)!على مسار} هو
مجموعها، $W_{AB}(\vect F) = \sum \scal{F}{\dd\ell}$ --- وهو ينهار على
\omterm{def:g10:relative-motion:trajectory}{مسار} مستقيم \omterm{def:g10:inertia:force}{بقوة} ثابتة إلى صيغة السنة الماضية
$F \times AB \times \cos\theta$.
\end{definition}

\begin{omfigure}
\begin{tikzpicture}[font=\small]
  \draw[omExo!60, thick] plot [smooth, tension=0.8] coordinates {(0,0.3) (1.2,1.15) (2.6,1.3) (4.0,0.7) (5.3,1.05)};
  \draw[omMeth, thick] (0,0.3) -- (1.2,1.15) -- (2.6,1.3) -- (4.0,0.7) -- (5.3,1.05);
  \foreach \p in {(0,0.3),(1.2,1.15),(2.6,1.3),(4.0,0.7),(5.3,1.05)} \fill \p circle (1.3pt);
  \node[below] at (0,0.3) {$A$}; \node[below right] at (5.3,1.05) {$B$};
  \draw[-Stealth, omThm, very thick] (2.6,1.3) -- (3.44,0.94) node[below left=-2pt] {$\vect{\dd\ell}$};
  \draw[-Stealth, omDef, very thick] (2.6,1.3) -- ++(22:1.5) node[above] {$\vect F$};
  \draw (2.6,1.3) ++(-23:0.75) arc[start angle=-23, end angle=22, radius=0.75];
  \node at ($(2.6,1.3)+(0:1.0)$) {$\theta$};
\end{tikzpicture}

{\small \omterm{def:g10:relative-motion:trajectory}{المسار} المنحني خط مكسور من خطوات قصيرة كثيرة: فتكسب كل
$\vect{\dd\ell}$ المقدارَ $F\,\dd\ell\cos\theta$؛ \omterm{def:g11:work-of-force:work}{والشغل} هو المجموع.}
\end{omfigure}

\begin{proposition}[شغل الوزن على أي مسار]\label{prop:g12:work-and-energy:weight}
على \emph{أي} \omterm{def:g10:relative-motion:trajectory}{مسار} من $A$ (على الارتفاع $z_A$) إلى $B$ (على الارتفاع
$z_B$)، ينجز \omterm{def:g10:universal-gravitation:weight}{وزن} كتلة $m$ \omterm{def:g11:work-of-force:work}{الشغل}
$W_{AB}(\vect P) = mg\,(z_A - z_B)$: فيدخل الهبوط، ولا يدخل الطريق أبدًا.
\end{proposition}

\begin{proof}
تسهم كل خطوة قصيرة بالمقدار $mg \times (\text{الهبوط الصغير})$ (وهو حساب
القطعة المستقيمة في السنة الماضية)؛ وتتداخل الهبوطات لتعطي
$z_A - z_B$. فالحالة المنحنية، المقبولة عندئذٍ، صارت مسلَّمة.
\end{proof}

\begin{definition}[القوى المحافظة وغير المحافظة]\label{def:g12:work-and-energy:conservative}
تكون \omterm{def:g10:inertia:force}{القوة} \emph{محافظة}\index{قوة محافظة} إذا كان شغلها
بين نقطتين واحدًا على كل \omterm{def:g10:relative-motion:trajectory}{مسار} (ومعدومًا على كل رحلة ذهاب
وإياب)، و\emph{غير محافظة}\index{قوة غير محافظة}
فيما عدا ذلك. \omterm{def:g10:universal-gravitation:weight}{فالوزن} \omterm{def:g10:inertia:force}{وقوة} الزنبرك (في القسم التالي): محافظتان.
أما \omterm{def:g10:inertia:inventory}{احتكاك} الانزلاق فلا: إذ يعاكس الحركة، فينجز $W = -fL$ على
\omterm{def:g10:relative-motion:trajectory}{مسار} طوله $L$ --- أي رسمًا على كل \omterm{def:g10:orders-of-magnitude:unit}{متر}، لا يُردّ أبدًا.
\end{definition}

\begin{example}[دربان إلى الكوخ]\label{ex:g12:work-and-energy:trails}
متنزّه كتلته \qty{75}{kg} يكسب \qty{300}{m} على درب حادّ طوله \qty{800}{m}
أو \qty{2.4}{km} من المنعطفات المتعرّجة. ويفرض \omterm{def:g10:universal-gravitation:weight}{الوزن}
$-75 \times 9.81 \times 300 \approx \qty{-2.2e5}{J}$ على كليهما؛ \omterm{def:g10:inertia:force}{وقوة} \omterm{def:g10:inertia:inventory}{احتكاك}
قدرها \qty{30}{N} تفرض \qty{-2.4e4}{J} على أحدهما و\qty{-7.2e4}{J}
على الآخر. \omterm{def:g10:universal-gravitation:weight}{والوزن} يردّ في النزول؛ أما فاتورة \omterm{def:g10:inertia:inventory}{الاحتكاك} فتذهب
حرارةً، في الاتجاهين.
\end{example}

\section{الطاقة الكامنة}

\begin{definition}[الطاقة الكامنة]\label{def:g12:work-and-energy:ep}
تُرفَق بكل \omterm{def:g12:work-and-energy:conservative}{قوة محافظة}
\emph{طاقة كامنة}\index{طاقة كامنة!من شغل محافظ}
$E_p$، وهي عدد لا يتوقف إلا على الموضع، ويُعرَّف بالعلاقة
$W_{AB} = E_p(A) - E_p(B)$: \omterm{def:g11:work-of-force:work}{فالشغل} الموصَّل \omterm{def:g10:energy-conservation:energy}{طاقة} كامنة منفَقة.
وتثبيت $E_p = 0$ عند نقطة مرجعية مختارة يثبّتها في كل مكان؛ و
الاختيار يزيح كل القيم بثابت يتلاشى من كل
فرق. وهذا يعيد في حالة \omterm{def:g10:universal-gravitation:weight}{الوزن} صيغة السنة الماضية $E_p = mgz$،
صحيحةً الآن على أي \omterm{def:g10:relative-motion:trajectory}{مسار}. ولا يمكن أن تكون \omterm{def:g10:inertia:inventory}{للاحتكاك} $E_p$: فالفاتورة
المتوقفة على الطريق ليست فرقًا بين عددين عند الطرفين --- ولا تخزّن
بقابلية الاسترجاع إلا القوى \omterm{def:g12:work-and-energy:conservative}{المحافظة}، وهو بالضبط ما يحافظ عليه الاسم.
\end{definition}

\begin{proposition}[الطاقة الكامنة المرنة]\label{prop:g12:work-and-energy:elastic}
زنبرك \omterm{def:g12:oscillators-and-time:spring}{صلابته} $k$ (\unit{N/m}) ممدود أو مضغوط بمقدار $x$
يسحب عائدًا \omterm{def:g10:inertia:force}{بالقوة} $F = kx$
(\cref{ch:g12:oscillators-and-time})؛ \omterm{def:g11:work-of-force:work}{والشغل} الذي يوصّله عائدًا
إلى السكون --- أي
\emph{طاقته الكامنة المرنة}\index{طاقة كامنة!مرنة} --- هو
$E_p = \tfrac12\, k x^2$.
\end{proposition}

\begin{proof}
تؤثر \omterm{def:g12:oscillators-and-time:pendulum}{القوة المرجِعة} على امتداد الحركة، ومن ثم يكون $\sum F\,\dd\ell$
هو \emph{المساحة تحت مخطط $F$--$x$}: أي مثلث قاعدته $x$ و
ارتفاعه $kx$، ومساحته $\tfrac12 x \times kx$. \omterm{def:g10:inertia:force}{والقوة} عند كل استطالة
هي نفسها ذهابًا وإيابًا: \omterm{def:g11:work-of-force:work}{فالشغل} تثبّته الطرفان، \omterm{def:g10:inertia:force}{وقوة} الزنبرك
\omterm{def:g12:work-and-energy:conservative}{محافظة}، و$E_p$ معرَّفة تعريفًا جيدًا.
\end{proof}

\begin{omfigure}
\begin{tikzpicture}
\begin{axis}[omaxis, width=8.4cm, height=4.8cm, xmin=0, xmax=0.095,
    ymin=0, ymax=21,
    xlabel={الاستطالة $x$ (\unit{m})}, ylabel={القوة $F$ (\unit{N})}]
  \addplot[draw=none, fill=omDef!18, domain=0:0.06, samples=2] {200*x} \closedcycle;
  \addplot[omThm, thick, domain=0:0.09, samples=2] {200*x};
  \draw[dashed, omExo] (axis cs:0,12) -- (axis cs:0.06,12) -- (axis cs:0.06,0);
  \node[font=\small, anchor=west] at (axis cs:0.061,6.5) {$kx$};
  \node[font=\small] at (axis cs:0.044,3.2) {$\tfrac12 kx^2$};
  \node[font=\small, anchor=south, omThm] at (axis cs:0.082,17) {$F = kx$};
\end{axis}
\end{tikzpicture}

{\small \omterm{def:g10:inertia:force}{قوة} الزنبرك تنمو مع الاستطالة؛ \omterm{def:g10:energy-conservation:energy}{والطاقة} المخزونة
هي المساحة تحت المستقيم: أي مثلث، $\tfrac12 \times x \times kx$.}
\end{omfigure}

\section{مبرهنتا الطاقة}

\begin{theorem}[مبرهنة الطاقة الحركية]\label{thm:g12:work-and-energy:ketheorem}\index{مبرهنة الطاقة الحركية}
في أي حركة لجسم كتلته $m$ من $A$ إلى $B$ --- منحنيةً كانت أو
ملتفّة أو ثلاثية الأبعاد ---
\[
\Delta E_k = E_k(B) - E_k(A) = \sum W_{AB}(\vect F),
\]
والمجموع على كل القوى المؤثرة في الجسم.
\end{theorem}

\begin{proof}
بما أن $v^2 = \scal{v}{v}$، تعطي قاعدة الجداء في رياضيات هذه السنة
مشتقةَ $v^2$ بالصورة $2\,\scal{a}{v}$، ومنه
$\dd E_k/\dd t = m\,\scal{a}{v} = \scal{F_{\text{الكلية}}}{v}$ \omterm{thm:g12:newtons-laws:second}{بقانون
نيوتن الثاني} (\cref{ch:g12:newtons-laws}). وفي زمن قصير $\dd t$
ينتقل الجسم $\vect{\dd\ell} = \vect v\,\dd t$: فتكسب $E_k$
\omterm{def:g11:work-of-force:work}{الشغل} الصغير \omterm{def:g10:inertia:force}{للقوة} \cref{def:g12:work-and-energy:path}؛ وجمع الخطوات
يجمع الأشغال --- وهي عبارة السنة الماضية المقبولة، مكسوبةً كاملة.
\end{proof}

\begin{proposition}[الاستطاعة اللحظية]\label{prop:g12:work-and-energy:power}
في كل لحظة توصّل \omterm{def:g10:inertia:force}{قوة} $\vect F$ على جسم سرعته $\vect v$
الاستطاعةَ $P = \scal{F}{v} = Fv\cos\theta$، و
$\dd E_k/\dd t$ هو مجموع استطاعات كل القوى: أي صيغة السنة الماضية
عند \omterm{def:g12:kinematics-2d:velocity}{السرعة} الثابتة، مضبوطةً في كل لحظة من أي حركة.
\end{proposition}

\begin{proof}
اقرأها من البرهان السابق: فكل \omterm{def:g10:inertia:force}{قوة} تغذّي
$\scal{F}{\dd\ell} = \scal{F}{v}\,\dd t$ في $E_k$ خلال $\dd t$.
\end{proof}

\begin{theorem}[مبرهنة الطاقة الميكانيكية]\label{thm:g12:work-and-energy:emtheorem}\index{مبرهنة الطاقة الميكانيكية}
ليكن $E_m = E_k + E_p$، حيث يجمع $E_p$ الطاقات الكامنة
لكل القوى \omterm{def:g12:work-and-energy:conservative}{المحافظة} العاملة. عندئذٍ، على أي \omterm{def:g10:relative-motion:trajectory}{مسار}،
\[
\Delta E_m = W_{AB}(\text{القوى غير المحافظة}),
\]
وتنحفظ $E_m$ كلما لم تنجز القوى \omterm{def:g12:work-and-energy:conservative}{غير المحافظة} شغلًا.
\end{theorem}

\begin{proof}
جزّئ مبرهنة \omterm{def:g11:mechanical-energy:kinetic}{الطاقة الحركية}: $\Delta E_k = W_{\text{المحافظة}} +
W_{\text{غير المحافظة}} = -\Delta E_p + W_{\text{غير المحافظة}}$؛ وانقل $\Delta E_p$ إلى الطرف الآخر.
\end{proof}

\begin{example}[النوّاس، بأمانة هذه المرة]\label{ex:g12:work-and-energy:pendulum}
كرة نوّاس على سلك طوله $L = \qty{2.5}{m}$ تُطلق من
السكون عند $\qty{45}{\degree}$ عن الوضع الشاقولي. \omterm{def:g10:inertia:inventory}{والشدّ}،
وهو عمودي على الحركة، عديم \omterm{def:g10:energy-conservation:power}{الاستطاعة}؛ \omterm{def:g10:universal-gravitation:weight}{والوزن} محافظ:
فتنحفظ $E_m$ \emph{على القوس المنحني}، وهو ما لم يكن يمكن في السنة الماضية
إلا التحقق منه. وتبدأ الكرة على ارتفاع $h = L(1 - \cos\qty{45}{\degree}) =
\qty{0.73}{m}$، ومن ثم عند الأسفل $v = \sqrt{2gh} \approx
\qty{3.8}{m/s}$.
\end{example}

\begin{method}[تدقيق طاقي لأي حركة]\label{met:g12:work-and-energy:audit}
\begin{enumerate}
  \item اختر الجملة وأحصِ القوى المؤثرة فيها.
  \item صنّف: فالعمودية على الحركة $\to$ لا شغل لها؛ \omterm{def:g12:work-and-energy:conservative}{والمحافظة}
        $\to$ تعطي $E_p$ في $E_m$؛ والباقية $\to$ $W_{\text{غير المحافظة}}$
        (\omterm{def:g10:inertia:inventory}{واحتكاك} الانزلاق: $-fL$، حيث $L$ = طول \omterm{def:g10:relative-motion:trajectory}{المسار}).
  \item اكتب $\Delta E_m = W_{\text{غير المحافظة}}$ بين نقطة المعطيات و
        نقطة السؤال؛ وحُلّ. \omterm{def:g10:energy-conservation:energy}{فالطاقة} تجيب عن ``كم \omterm{def:g12:kinematics-2d:velocity}{السرعة} وكم
        الارتفاع''؛ ولا يحتاج قوانين نيوتن إلا الزمن والتسارع.
\end{enumerate}
\end{method}

\section{مخططات الطاقة}

\begin{definition}[مخطط الطاقة ونقاط الارتداد]\label{def:g12:work-and-energy:diagram}
في حركة على محور واحد $x$ مع $E_m$ محفوظة، يرسم
\emph{مخطط الطاقة}\index{مخطط الطاقة} المنحني $E_p(x)$
والخط الأفقي $E_m$. وبما أن $E_k = E_m - E_p \geq 0$، تنحصر
الحركة حيث يقع المنحني تحت الخط؛ والفجوة
بينهما هي $E_k$، مقروءةً مباشرة. وعند
\emph{نقطة ارتداد}\index{نقطة ارتداد}، حيث يلاقي المنحني الخط،
تكون $v = 0$ وترتدّ الحركة.
\end{definition}

\begin{omfigure}
\begin{tikzpicture}
\begin{axis}[omaxis, width=9.0cm, height=5.4cm,
    xlabel={$x$}, ylabel={الطاقة}, xtick=\empty, ytick=\empty,
    xmin=-2.25, xmax=2.25, ymin=0, ymax=2.45]
  \addplot[omDef, thick, domain=-2.05:2.05, samples=80] {0.25*x^4 - x^2 + 1.5};
  \addplot[omThm, dashed, thick, domain=-2.25:2.25, samples=2] {1.0};
  \addplot[omProp, dashed, domain=-2.25:2.25, samples=2] {1.7};
  \node[font=\small, omThm, anchor=west] at (axis cs:-2.2,0.86) {$E_m$};
  \node[font=\small, omProp, anchor=west] at (axis cs:-2.2,1.85) {$E_m'$};
  \node[font=\small, omDef, anchor=west] at (axis cs:1.55,2.1) {$E_p(x)$};
  \fill (axis cs:0.765,1.0) circle (1.6pt) node[font=\small, below=1pt] {$x_1$};
  \fill (axis cs:1.848,1.0) circle (1.6pt) node[font=\small, below right=-1pt] {$x_2$};
  \fill[omMeth] (axis cs:1.414,0.5) circle (2.2pt) node[font=\small, below] {مستقر};
  \fill[omMeth] (axis cs:0,1.5) circle (2.2pt) node[font=\small, above] {غير مستقر};
\end{axis}
\end{tikzpicture}

{\small مخطط \omterm{def:g10:energy-conservation:energy}{طاقة}، مقروء دون حلّ أي شيء: فبالطاقة
$E_m$ يتذبذب الجسم بين $x_1$ و$x_2$، وأسرع ما يكون حيث تكون
الفجوة أكبر؛ وبالطاقة $E_m'$ يزحف فوق التلّ \omterm{rem:g12:work-and-energy:escape}{ويفلت}.}
\end{omfigure}

\begin{proposition}[القوة من المنظر؛ أوضاع التوازن]\label{prop:g12:work-and-energy:equilibria}
\omterm{def:g12:work-and-energy:conservative}{القوة المحافظة} على امتداد $x$ هي سالب ميل $E_p$:
أي $F = -\,\dd E_p/\dd x$ --- فهي تتجه \emph{نازلةً} على المخطط،
وتنعدم حيث يكون المماسّ أفقيًّا: أي عند توازن. والصغرى
في $E_p$ توازن \emph{مستقر}\index{توازن مستقر} --- إذ يُدفع
الجسم المزاح عائدًا، ككرة في وادٍ؛ والعظمى توازن
\emph{غير مستقر}\index{توازن غير مستقر} --- إذ يُدفع بعيدًا، ككرة
على قمة تلّ.
\end{proposition}

\begin{proof}
على خطوة صغيرة $\dd x$ تنجز \omterm{def:g10:inertia:force}{القوة} \omterm{def:g11:work-of-force:work}{الشغل} $F\,\dd x = E_p(x) -
E_p(x + \dd x) = -\dd E_p$؛ فاقسم على $\dd x$. والنزول على أي من جانبَي
وادٍ يشير عائدًا إلى الداخل؛ وحول قمة تلّ، إلى الخارج.
\end{proof}

\begin{remark}[التذبذبات والإفلات]\label{rem:g12:work-and-energy:escape}
التذبذبات الصغيرة في \cref{ch:g12:oscillators-and-time} تسكن
في قيعان وديان $E_p$. وبعيدًا عن الأرض تضعف \omterm{def:g11:fundamental-interactions:four}{الجاذبية}: فيصعد
منحني $E_p(r)$ ببطء متزايد نحو سقف منتهٍ (فالمساحة
تحت مخطط \omterm{def:g10:inertia:force}{القوة} المتقلص منتهية). والصاروخ الذي تبلغ $E_m$ الخاصة به
السقفَ لا يلاقي \omterm{def:g12:work-and-energy:diagram}{نقطة ارتداد} أبدًا: فهو
\emph{يفلت}\index{سرعة الإفلات} --- بنحو \qty{11.2}{km/s} من
الأرض (\cref{exo:g12:work-and-energy:13})؛ ودون ذلك يبقى
مقيَّدًا. أما السقف نفسه فيُحسب في مجلد السنة الأولى.
\end{remark}

\begin{example}[الحلقة الرأسية]\label{ex:g12:work-and-energy:loop}
كرة زجاجية تُطلق من الارتفاع $H$ تنساب إلى حلقة عمودية نصف
قطرها $R$. ويقتضي ديناميك الدوران (\cref{ch:g12:newtons-laws})
أن $v_{\text{القمة}}^2 \geq gR$ عند القمة --- وهو المعطى الذي استعارته مسألة
الأفعوانية في السنة الماضية، وقد صار لنا الآن. ويعطي الانحفاظ من السكون عند $H$
إلى القمة (على الارتفاع $2R$) العلاقةَ $v_{\text{القمة}}^2 = 2g(H - 2R)$، ومنه
$H \geq \tfrac52 R$: أي نصف نصف قطر فوق قمة الحلقة، مهما تكن
الكتلة --- وأكثر قليلًا في العالم الحقيقي المحتكّ.
\end{example}

\begin{omfigure}
\begin{tikzpicture}[font=\small, scale=0.9]
  \draw[thick] (0,3.0) .. controls (0.9,1.1) and (1.6,0) .. (3.1,0) -- (6.9,0);
  \draw[thick] (5.5,1.2) circle (1.2);
  \fill[omDef] (0,3.0) circle (2.2pt);
  \draw[dashed, omExo] (4.0,2.4) -- (7.6,2.4) node[right] {$2R$};
  \draw[dashed, omThm] (-0.4,3.0) -- (7.6,3.0) node[right] {$H_{\min} = \tfrac52 R$};
  \draw[Stealth-Stealth, omThm] (7.35,2.4) -- (7.35,3.0);
  \node[right, omThm] at (7.35,2.7) {$\tfrac{R}{2}$};
  \draw[-Stealth, omMeth, thick] (5.5,2.4) -- (4.5,2.4) node[above right=0pt and 2pt] {$v_{\text{القمة}}$};
  \fill (5.5,2.4) circle (1.5pt);
\end{tikzpicture}

{\small الحلقة الرأسية: تقع نقطة الإطلاق على نصف نصف قطر فوق القمة،
فتُبقي $E_k$ الفائضة الكرةَ مضغوطة ($v_{\text{القمة}}^2 = gR$).}
\end{omfigure}

\section{تمارين}

\begin{exercise}[$\star$]\label{exo:g12:work-and-energy:1}
زنبرك \omterm{def:g12:oscillators-and-time:spring}{صلابته} $k = \qty{200}{N/m}$ ممدود بمقدار
\qty{5.0}{cm}. احسب \omterm{def:g10:energy-conservation:energy}{الطاقة} المخزونة؛ وماذا تصير إذا تضاعفت
الاستطالة؟ وكم تكلّف \emph{الاستطالة} \qty{5.0}{cm} \omterm{def:g10:orders-of-magnitude:unit}{الثانية}؟
\end{exercise}

\begin{exercise}[$\star$]\label{exo:g12:work-and-energy:2}
مسيّرة كتلتها \qty{0.90}{kg} تطير على \omterm{def:g10:relative-motion:trajectory}{مسار} متعرّج من شرفة على
$z = \qty{12}{m}$ إلى شرفة على $z = \qty{30}{m}$. فما شغل
وزنها؟ وعلى الصعود المستقيم؟ وفي رحلة الذهاب والإياب؟
\end{exercise}

\begin{exercise}[$\star$]\label{exo:g12:work-and-energy:3}
\omterm{def:g12:work-and-energy:conservative}{محافظة} أم \omterm{def:g12:work-and-energy:conservative}{غير محافظة} أم عديمة \omterm{def:g11:work-of-force:work}{الشغل}؟ برّر من المسارات:
(أ) \omterm{def:g10:universal-gravitation:weight}{الوزن}؛ (ب) \omterm{def:g10:inertia:force}{قوة} الزنبرك؛ (ج) \omterm{def:g10:inertia:inventory}{احتكاك} الانزلاق؛ (د)
\omterm{def:g12:newtons-laws:friction}{القوة الناظمية} على صندوق منزلق؛ (ه) \omterm{def:g11:circuits-and-power:resistance}{مقاومة} الهواء.
\end{exercise}

\begin{exercise}[$\star$]\label{exo:g12:work-and-energy:4}
نوّاس طوله \qty{2.0}{m} يُطلق من السكون عند
$\qty{60}{\degree}$ عن الوضع الشاقولي. جِد سرعته عند أخفض
نقطة. ولماذا لا يدخل \omterm{def:g10:inertia:inventory}{شدّ} السلك في الميزانية أبدًا؟
\end{exercise}

\begin{exercise}[$\star$]\label{exo:g12:work-and-energy:5}
تركب راكبة دراجة \omterm{def:g12:kinematics-2d:velocity}{بسرعة} \qty{9.0}{m/s} على المستوي، وتوصّل
\qty{250}{W}: فأيّ \omterm{def:g10:inertia:force}{قوة} \omterm{def:g11:circuits-and-power:resistance}{مقاومة} كلية تصارع؟ وإذا صعدت \omterm{def:g12:kinematics-2d:velocity}{بسرعة}
\qty{5.0}{m/s} \omterm{def:g10:energy-conservation:power}{بالاستطاعة} نفسها: فأيّ \omterm{def:g10:inertia:force}{قوة} تتغلب عليها؟
\end{exercise}

\begin{exercise}[$\star\star$]\label{exo:g12:work-and-energy:6}
زنبرك قاذفة لعبة ($k = \qty{450}{N/m}$) يُضغط بمقدار
\qty{8.0}{cm} خلف سهم كتلته \qty{20}{g}. جِد \omterm{def:g10:energy-conservation:energy}{الطاقة} المخزونة، \omterm{def:g12:kinematics-2d:velocity}{وسرعة}
الإطلاق، والارتفاع المبلوغ إذا قُذف شاقوليًّا إلى أعلى (بلا هواء).
\end{exercise}

\begin{exercise}[$\star\star$]\label{exo:g12:work-and-energy:7}
بمطّ شريط مطاطي، يقيس تلميذ $F = 0$، $3.0$، $7.0$،
$12.0$، \qty{18.0}{N} عند $x = 0$، $2.0$، $4.0$، $6.0$، \qty{8.0}{cm}.
قدّر \omterm{def:g10:energy-conservation:energy}{الطاقة} المخزونة بمساحات أشباه المنحرفات؛ وقارنها بالمقدار
$\tfrac12 kx^2$، بأخذ $k = F/x$ عند النهاية. ولماذا تختلفان؟
\end{exercise}

\begin{exercise}[$\star\star$]\label{exo:g12:work-and-energy:8}
متزلج بزحّافة كتلته \qty{55}{kg} يبدأ ساكنًا، ويهبط \qty{8.0}{m} على
\omterm{def:g10:relative-motion:trajectory}{مسار} منحنٍ طوله \qty{25}{m}، فيصل \omterm{def:g12:kinematics-2d:velocity}{بسرعة} \qty{9.0}{m/s}. فما \omterm{def:g11:mechanical-energy:em}{الطاقة
الميكانيكية} المفقودة؟ وما \omterm{def:g10:inertia:force}{قوة} \omterm{def:g10:inertia:inventory}{الاحتكاك} المتوسطة؟ ولماذا لا يهمّ شكل \omterm{def:g10:relative-motion:trajectory}{المسار} بالضبط؟
\end{exercise}

\begin{exercise}[$\star\star$]\label{exo:g12:work-and-energy:9}
لمضمار كرات زجاجية حلقة نصف قطرها \qty{20}{cm}. جِد أدنى
ارتفاع إطلاق (من السكون، بلا \omterm{def:g10:inertia:inventory}{احتكاك}). وإذا أُطلقت من \qty{60}{cm}:
فما \omterm{def:g12:kinematics-2d:velocity}{السرعة} عند قمة الحلقة، وما النسبة $v_{\text{القمة}}^2/(gR)$؟ ولماذا يجب
أن يُطلق مضمار حقيقي من ارتفاع أعلى بعدُ؟
\end{exercise}

\begin{exercise}[$\star\star$]\label{exo:g12:work-and-energy:10}
خرزة كتلتها \qty{100}{g} تنزلق بلا \omterm{def:g10:inertia:inventory}{احتكاك} على محور، حيث
$E_p(x) = 8.0\,x^2$ (\omterm{def:g11:work-of-force:work}{بالجول}، و$x$ \omterm{def:g10:orders-of-magnitude:unit}{بالمتر}) و$E_m = \qty{2.0}{J}$.
فما نقاط الارتداد؟ وما \omterm{def:g12:kinematics-2d:velocity}{السرعة} القصوى؟ وأيّ نوع من الحركات هذا؟
\end{exercise}

\begin{exercise}[$\star\star$]\label{exo:g12:work-and-energy:11}
سيارة كتلتها \qty{1200}{kg} \omterm{def:g12:kinematics-2d:velocity}{بسرعة} \qty{2.0}{m/s} يوقفها زنبرك مصدّ
\omterm{def:g12:oscillators-and-time:spring}{صلابته} $k = \qty{6.0e5}{N/m}$. جِد أقصى انضغاط.
وعند \qty{4.0}{m/s}: لماذا ليس أربعة أضعاف، مع أن \omterm{def:g10:energy-conservation:energy}{الطاقة} تتربّع؟
\end{exercise}

\begin{exercise}[$\star\star\star$]\label{exo:g12:work-and-energy:12}
عربة كتلتها \qty{0.50}{kg} تتدحرج على مضمار له $E_p(x)$ فيه
قمة تلّ \qty{6.0}{J} ($x = 0$) ووادٍ \qty{2.0}{J}
($x = \qty{1.5}{m}$). (أ) فإذا وصلت من اليسار \omterm{def:g10:energy-conservation:energy}{بطاقة}
$E_m = \qty{5.0}{J}$: هل تعبر؟ وماذا يحدث بدل ذلك؟ (ب) وبطاقة
$E_m = \qty{7.0}{J}$: ما السرعتان عند قمة التلّ وفي الوادي؟ (ج) وما أوضاع التوازن؟
\end{exercise}

\begin{exercise}[$\star\star\star$]\label{exo:g12:work-and-energy:13}
بقبول أن جرّ كتلة $m$ من الأرض إلى بُعد كيفي
يكلّف $mgR_E$، $R_E = \qty{6.37e6}{m}$ (وهي المساحة المنتهية تحت مخطط
\omterm{def:g11:fundamental-interactions:four}{الجاذبية} المتضائلة)، احسب \emph{\omterm{rem:g12:work-and-energy:escape}{سرعة الإفلات}} من الأرض.
ولماذا هي نفسها للمسبار وللحصاة؟ وماذا عن إطلاق دونها بقليل؟
\end{exercise}

\begin{exercise}[$\star\star\star$]\label{exo:g12:work-and-energy:14}
منزلقة كتلتها \qty{250}{g} تتذبذب على زنبرك ($k = \qty{40}{N/m}$)
\omterm{def:g12:oscillators-and-time:oscillator}{بسعة} \qty{6.0}{cm} (\cref{ch:g12:oscillators-and-time}).
فما \omterm{def:g10:energy-conservation:energy}{الطاقة} الكلية؟ وما \omterm{def:g12:kinematics-2d:velocity}{السرعة} القصوى؟ وما المواضع التي عندها $E_k = E_p$؟
\end{exercise}

\begin{exercise}[$\star\star\star$]\label{exo:g12:work-and-energy:15}
يعدو قافز بالزانة \omterm{def:g12:kinematics-2d:velocity}{بسرعة} \qty{9.5}{m/s}. قدّر الارتفاع الذي يكسبه
مركز كتلته إذا حوّلت الزانة كل طاقته الحركية، و
العارضة المتجاوَزة (\omterm{def:g11:forces-and-motion:com}{ومركز الكتلة} يبدأ على ارتفاع \qty{1.0}{m}). والرقم القياسي
نحو \qty{6.3}{m}: فما الذي يلتقطه التقدير، وما الذي يهمله؟
\end{exercise}

\section{مسألة: قفزة التزلج}

\begin{problem}[{مسألة نهاية الأسبوع --- قفزة التزلج: مسار اندفاع يدقّقه
مسدس سرعة، وطيران في سقوط حرّ، وهندسة لطيفة لهبوط حادّ،
والبرج الذي يجب أن يشتريه المصمّم متى فوتر الاحتكاك}]
\label{pb:g12:work-and-energy:1}
قافز تزلج أولمبي ($m = \qty{70}{kg}$، مع الزلاجتين) يبدأ
ساكنًا من بوابة على ارتفاع $H = \qty{45}{m}$ فوق طاولة الانطلاق، وينزلق
$L = \qty{90}{m}$ من \omterm{def:g10:relative-motion:trajectory}{مسار} الاندفاع المنحني، ويغادر الطاولة أفقيًّا؛
ويقرأ مسدس \omterm{def:g12:kinematics-2d:velocity}{سرعة} هناك $v_0 = \qty{26}{m/s}$. وتقع نقطة الهبوط على انخفاض
$h = \qty{40}{m}$، على منحدر ميله $\qty{35}{\degree}$. ويهمل الجزآن II
وIII \omterm{def:g11:circuits-and-power:resistance}{مقاومة} الهواء.

\medskip\noindent\textbf{الجزء I --- \omterm{def:g10:relative-motion:trajectory}{مسار} الاندفاع.}
\begin{enumerate}
  \item تنبّأ \omterm{def:g12:kinematics-2d:velocity}{بسرعة} الانطلاق بلا \omterm{def:g10:inertia:inventory}{احتكاك}؛ ولماذا لا يهمّ منحني
        \omterm{def:g10:relative-motion:trajectory}{مسار} الاندفاع؟
  \item \omterm{def:g10:relative-motion:frame}{بالمرجع} عند الطاولة: احسب $E_m$ عند البوابة وعند الانطلاق.
  \item احسب $\Delta E_m$؛ وأيّ مبرهنة تسمّي الجناة، ومن هم؟
  \item استنتج \omterm{def:g10:inertia:force}{قوة} \omterm{def:g11:circuits-and-power:resistance}{المقاومة} المتوسطة (الثلج مع الهواء) على \omterm{def:g10:relative-motion:trajectory}{المسار}.
  \item وأيّ جزء من $E_m$ يضيع؟ ولماذا نشمّع الزلاجتين وننحني؟
\end{enumerate}

\medskip\noindent\textbf{الجزء II --- الطيران.}
\begin{enumerate}[resume]
  \item لا يؤثر في الطيران إلا \omterm{def:g10:universal-gravitation:weight}{الوزن}: فأيّ حركة من حركات
        \cref{ch:g12:projectile-motion} هذه؟ وما مركّبتا \omterm{def:g12:kinematics-2d:velocity}{السرعة}
        عند الهبوط؟
  \item ما \omterm{def:g12:kinematics-2d:velocity}{سرعة} الهبوط، مرتين: من المركّبتين، ومن انحفاظ
        $E_m$؟ وبوحدة \unit{km/h}؟
  \item احسب \omterm{def:g12:projectile-motion:range}{زمن الطيران} والمسافة الأفقية المقطوعة.
  \item ويركب القافزون الحقيقيون الهواءَ كجناح، فيطيرون أبعد بكثير:
        فهل يقدر الهواء أن \emph{يرفع} \omterm{def:g12:kinematics-2d:velocity}{سرعة} الهبوط؟ ناقش
        بمبرهنة \omterm{def:g11:mechanical-energy:em}{الطاقة الميكانيكية}.
\end{enumerate}

\medskip\noindent\textbf{الجزء III --- الهبوط.}
\begin{enumerate}[resume]
  \item احسب زاوية \omterm{def:g12:kinematics-2d:velocity}{سرعة} الهبوط تحت الأفقي.
  \item فكّكها إلى مركّبتين موازية وعمودية على
        المنحدر.
  \item ولا تمتصّ الساقان إلا الجزء العمودي: احسب تلك
        \omterm{def:g11:mechanical-energy:kinetic}{الطاقة الحركية}؛ وقارنها بالكلية.
  \item أما الأرض المستوية فتمتصّ كل شيء: فأيّ هبوط شاقولي يعادل
        كل هبوط ($h_{\text{المكافئ}} = v_\perp^2/2g$،
        على الترتيب\ $v^2/2g$)؟
  \item وفي جملة واحدة: لماذا تكون تلال الهبوط حادّة ومنحنية؟
\end{enumerate}

\medskip\noindent\textbf{الجزء IV --- برج المصمّم.}
تلّ أصغر يحتاج انطلاقًا عند $v_0' = \qty{24}{m/s}$ لا غير؛ وينزل
\omterm{def:g10:relative-motion:trajectory}{مسار} اندفاعه المستقيم بميل $\qty{35}{\degree}$، ومن ثم يعني ارتفاع البوابة $H'$
طول مضمار $H'/\sin\qty{35}{\degree}$؛ وبالمقاومة \qty{80}{N} نفسها.
\begin{enumerate}[resume]
  \item احسب ارتفاع البوابة الذي سيبنيه مصمّم يهمل \omterm{def:g10:inertia:inventory}{الاحتكاك}.
  \item اكتب مبرهنة \omterm{def:g11:mechanical-energy:em}{الطاقة الميكانيكية}، من البوابة إلى الطاولة، من أجل $H'$.
  \item حُلّ لإيجاد $H'$.
  \item وبأي نسبة مئوية رفع \omterm{def:g10:inertia:inventory}{الاحتكاك} البرجَ؟
  \item ولم تعد الكتلة تتلاشى: أعِد حساب $H'$ من أجل
        قافز كتلته \qty{60}{kg}. فمن يحتاج البرج الأعلى، ولماذا؟
  \item بيت القصيد: البرج المعلَن، وحصة \omterm{def:g10:inertia:inventory}{الاحتكاك} منه.
\end{enumerate}
\end{problem}
